\documentclass[11pt,a4paper]{article}
\usepackage{fontspec}
\setmainfont{DejaVu Serif}
\setsansfont{DejaVu Sans}
\setmonofont{DejaVu Sans Mono}
\usepackage{polyglossia}
\setmainlanguage{polish}
\usepackage[a4paper,top=2.4cm,bottom=2.6cm,left=2.4cm,right=2.4cm]{geometry}
\usepackage{amssymb}
\usepackage{xcolor}
\definecolor{navy}{RGB}{0,51,102}
\definecolor{lightnavy}{RGB}{232,238,246}
\definecolor{warnred}{RGB}{153,0,0}
\usepackage{titlesec}
\titleformat{\section}{\Large\bfseries\sffamily\color{navy}}{\thesection.}{0.6em}{}
\titleformat{\subsection}{\large\bfseries\sffamily\color{navy}}{\thesubsection.}{0.6em}{}
\titleformat{\subsubsection}{\normalsize\bfseries\sffamily\color{navy}}{\thesubsubsection.}{0.6em}{}
\usepackage{enumitem}
\setlist{itemsep=2pt,topsep=4pt}
\usepackage{booktabs}
\usepackage{longtable}
\usepackage{array}
\usepackage{listings}
\lstset{basicstyle=\ttfamily\footnotesize,breaklines=true,frame=single,framerule=0.3pt,rulecolor=\color{navy},backgroundcolor=\color{lightnavy},tabsize=2,columns=fullflexible,keepspaces=true}
\usepackage{fancyhdr}
\pagestyle{fancy}
\fancyhf{}
\fancyhead[L]{\sffamily\small\textcolor{navy}{AMPERE POINT --- Przewodnik instalacji i integracji HA}}
\fancyhead[R]{\sffamily\small\textcolor{navy}{v1 --- 2026-07-06}}
\fancyfoot[C]{\sffamily\small\thepage}
\renewcommand{\headrulewidth}{0.4pt}
\usepackage{hyperref}
\hypersetup{colorlinks=true,linkcolor=navy,urlcolor=navy}
\setlength{\parskip}{4pt}
\setlength{\parindent}{0pt}
\emergencystretch=3em
\sloppy

\newcommand{\uwaga}[1]{\par\medskip\noindent\textcolor{warnred}{\textbf{UWAGA:}} #1\par\medskip}
\newcommand{\wskazowka}[1]{\par\medskip\noindent\textcolor{navy}{\textbf{Wskazówka:}} #1\par\medskip}

\begin{document}

\begin{center}
{\huge\bfseries\sffamily\color{navy} Home Assistant dla AMPERE POINT\\[2pt] --- kompletny przewodnik instalacji i integracji}\\[8pt]
{\large Od czystego Windowsa do działającego systemu: Docker, Home Assistant, HACS,\\ Tuya, xtend\_tuya, dodatek firmowy, panel Energia i automatyzacje\\[2pt] + rozszerzenia: kamera Imou i symulator falownika SUN2000}\\[6pt]
{\normalsize wersja 1 --- 6 lipca 2026 --- dokument wewnętrzny projektu AMPERE\_POINT}\\[2pt]
{\small Przewodnik pisany dla osoby bez wcześniejszego doświadczenia z Home Assistant i Dockerem.\\ Zawiera wszystkie pułapki znalezione w sesjach 2026-06-30 \ldots\ 2026-07-05 i sposoby ich ominięcia.}
\end{center}

\vspace{4pt}
\tableofcontents
\newpage

%=====================================================================
\section{Cel, zakres i jak korzystać z przewodnika}

Ten dokument prowadzi \textbf{od zera} --- od komputera z samym Windows --- do w pełni działającej integracji ładowarek AMPERE POINT (seria Q, moduły Tuya) z systemem automatyki \textbf{Home Assistant}: pełne dane (moce, prądy, energie, statusy), koszty ładowania w PLN, panel Energia i automatyzacje. Rozdziały opcjonalne opisują dołączenie kamery Imou oraz symulatora falownika Huawei SUN2000 (do ćwiczeń przed montażem fotowoltaiki).

Przewodnik powstał na bazie rzeczywistych wdrożeń w projekcie AMPERE\_POINT (sesje 2026-06-30 do 2026-07-05). Wszystkie opisane pułapki \textbf{wydarzyły się naprawdę} --- każdą oznaczono i podano sprawdzone obejście. Dokument uzupełnia (nie zastępuje) raport \texttt{AMPERE\_POINT\_home\_assistant\_xtend\_tuya\_raport\_v1}, który wyjaśnia szczegółowo, \emph{dlaczego} system działa tak, a nie inaczej.

\subsection{Konwencje zapisu}

\begin{itemize}
\item \textbf{Ustawienia $\to$ Urządzenia i usługi} --- ścieżka klikania w menu (kolejne pozycje po strzałce).
\item \texttt{czcionka maszynowa} --- adresy, polecenia, nazwy plików i katalogów; wpisywać dokładnie tak, jak wydrukowano.
\item Ramki z kodem --- polecenia do skopiowania w całości.
\item \textcolor{warnred}{\textbf{UWAGA}} --- pułapka potwierdzona w praktyce; przeczytać \emph{zanim} wykona się krok.
\item Terminy techniczne wyjaśniane są przy pierwszym użyciu i zebrane w słowniczku (rozdział~\ref{sec:slowniczek}).
\end{itemize}

\subsection{Czego przewodnik wymaga od czytelnika}

Umiejętności: obsługa przeglądarki, kopiowanie plików, uruchamianie programów. \textbf{Nic więcej.} Każde polecenie terminala jest podane w całości do skopiowania. Tam, gdzie system poprosi o uprawnienia administratora (okno UAC --- ciemniejszy ekran z pytaniem ,,Czy chcesz zezwolić\ldots''), wystarczy kliknąć \textbf{Tak}.

%=====================================================================
\section{Obraz całości --- co budujemy i dlaczego tak}
\label{sec:architektura}

\subsection{Elementy układanki}

\begin{description}[style=nextline]
\item[Home Assistant (HA)] darmowy system automatyki domowej: zbiera dane z urządzeń w jednym panelu www i pozwala nimi sterować oraz budować automatyzacje (,,gdy X, zrób Y'').
\item[Docker Desktop] program uruchamiający aplikacje w \textbf{kontenerach} --- odizolowanych ,,pudełkach'' z kompletem zależności. HA działa u nas jako kontener \texttt{homeassistant}; jego konfiguracja leży w zwykłym folderze \texttt{C:\textbackslash HomeAssistant\textbackslash config}.
\item[HACS] społecznościowy ,,sklep'' z dodatkami do HA, których nie ma w wersji oficjalnej. Stąd instalujemy xtend\_tuya (i w przyszłości Huawei Solar).
\item[Integracja Tuya (oficjalna)] łączy HA z chmurą Tuya, w której ,,mieszkają'' ładowarki serii Q. Sama z siebie daje \textbf{tylko 1 encję na urządzenie} (przełącznik) --- to normalne.
\item[xtend\_tuya] dodatek społecznościowy (autor azerty9971) dokładający brakujące encje: prądy, napięcia, moce, energie, statusy. W trybie pełnym korzysta z \textbf{platformy deweloperskiej Tuya} (darmowej).
\item[TuyaExtend AmperePoint] firmowy dodatek-nakładka: mapuje encje źródłowe na ,,ładną'' ładowarkę, liczy koszty sesji wg taryfy w PLN, tłumaczy statusy na polski.
\end{description}

Droga danych (szczegóły i protokoły --- raport v1, rozdz. 6): ładowarka $\to$ Wi-Fi 2,4 GHz $\to$ chmura Tuya (centrum danych Europa Środkowa) $\to$ internet $\to$ Home Assistant na komputerze. Integracja jest \textbf{chmurowa}: bez internetu encje będą ,,niedostępne'', choć ładowarki dalej ładują.

\subsection{Dlaczego Docker, a nie maszyna wirtualna}
\label{sec:dlaczego-docker}

Pierwsze podejście (2026-06-30) stawiało HA OS w maszynie wirtualnej VirtualBox. Instalacja \textbf{utknęła na wiele godzin} na pobieraniu Core: wirtualna karta sieciowa w trybie mostka na Wi-Fi dławiła transfer, a dodatkowo twarde wyłączenie maszyny (host zasnął) uszkodziło magazyn Supervisora. Wnioski, które przenieśliśmy do tego przewodnika:

\begin{itemize}
\item \textbf{Docker Desktop (WSL2) ma wydajną sieć} --- pobieranie obrazu HA trwa minuty, nie godziny. Wariant kontenerowy okazał się stabilny i to on jest ścieżką główną.
\item Podczas długich pobrań \textbf{komputer nie może zasnąć} (rozdz.~\ref{sec:etap0}).
\item Wersja kontenerowa nie ma sklepu dodatków HA OS ani przycisku ,,Uruchom ponownie'' w menu --- restart robi się w Docker Desktop (to normalne, nie awaria).
\end{itemize}

\subsection{Mapa etapów}

\begin{center}
\begin{tabular}{clc}
\toprule
\textbf{Etap} & \textbf{Co powstaje} & \textbf{Rozdział}\\
\midrule
0 & WSL2 + Docker Desktop & \ref{sec:etap0}\\
1 & Kontener Home Assistant + konto & \ref{sec:etap1}\\
2 & HACS & \ref{sec:etap2}\\
3 & Oficjalna integracja Tuya (urządzenia widoczne) & \ref{sec:etap3}\\
4 & xtend\_tuya bez chmury dev (encje $\times$15) & \ref{sec:etap4}\\
5 & Projekt na platformie deweloperskiej Tuya & \ref{sec:etap5}\\
6 & Tryb DP Instruction produktów & \ref{sec:etap6}\\
7 & Klucze OpenAPI w HA (encje ,,prawie wszystkie'') & \ref{sec:etap7}\\
8 & Dodatek firmowy TuyaExtend AmperePoint & \ref{sec:etap8}\\
9 & Porządki, panel Energia, pulpit & \ref{sec:etap9}\\
10 & Automatyzacje & \ref{sec:etap10}\\
--- & Opcjonalnie: kamera Imou, symulator SUN2000 & \ref{sec:kamera}, \ref{sec:pvsim}\\
\bottomrule
\end{tabular}
\end{center}

%=====================================================================
\section{Wymagania wstępne --- checklista przed startem}

\begin{itemize}
\item Komputer z \textbf{Windows 10/11 (64-bit)}, min. 8 GB RAM, ok. 10 GB wolnego dysku, \textbf{stały dostęp do internetu} (najlepiej kabel; Wi-Fi wystarczy dla Dockera).
\item Konto w aplikacji \textbf{Smart Life} (lub Imou/Tuya Smart) z dodanymi ładowarkami --- u nas konto \texttt{dawid.cekala@gmail.com} z 10 urządzeniami. Telefon z tą aplikacją pod ręką (będzie skanować kody QR).
\item Sieć Wi-Fi \textbf{2,4 GHz} dla urządzeń Tuya (moduły nie łączą się z 5 GHz).
\item Konto GitHub (darmowe) --- potrzebne raz, do autoryzacji HACS.
\item Ok. 2--3 godziny czasu na etapy 0--8 (większość to klikanie wg instrukcji).
\item Zasilanie: laptop podpięty do ładowarki; \textbf{uśpienie wyłączone} na czas instalacji (niżej).
\end{itemize}

%=====================================================================
\section{Etap 0 --- przygotowanie Windows: WSL2 i Docker Desktop}
\label{sec:etap0}

\subsection{Wyłącz usypianie komputera (na czas instalacji)}

\textbf{Ustawienia Windows $\to$ System $\to$ Zasilanie i bateria $\to$ Ekran i uśpienie} $\to$ pozycje ,,\emph{Przy zasilaniu\ldots{} uśpij}'' ustaw na \textbf{Nigdy}. Uśpienie komputera w trakcie pobierania było jedną z przyczyn porażki wariantu VM (rozdz.~\ref{sec:dlaczego-docker}). Po zakończeniu instalacji można przywrócić stare ustawienia --- ale jeśli HA ma działać całą dobę (zbierać historię, wykonywać automatyzacje), komputer musi pozostawać włączony na stałe.

\subsection{WSL2 --- silnik Linuksa dla Dockera}

WSL2 (Windows Subsystem for Linux) to wbudowany w Windows mechanizm uruchamiania lekkiego Linuksa; Docker Desktop używa go ,,pod maską''.

\begin{enumerate}
\item Kliknij \textbf{prawym} przyciskiem na przycisk \textbf{Start} $\to$ \textbf{Terminal (administrator)} (albo ,,Windows PowerShell (administrator)''). Potwierdź okno UAC przyciskiem \textbf{Tak}.
\item Wpisz i zatwierdź Enterem:
\begin{lstlisting}
wsl --install
\end{lstlisting}
\item Po zakończeniu \textbf{zrestartuj komputer}. Jeśli polecenie odpowie, że WSL już jest zainstalowany --- restart można pominąć.
\end{enumerate}

\subsection{Docker Desktop}

\begin{enumerate}
\item Pobierz instalator: \url{https://www.docker.com/products/docker-desktop/} (przycisk \textbf{Download for Windows}).
\item Uruchom instalator; zostaw zaznaczone \textbf{Use WSL 2 instead of Hyper-V}; potwierdź UAC. Po instalacji wyloguj się/zrestartuj, jeśli poprosi.
\item Uruchom \textbf{Docker Desktop} (ikona wieloryba). Przy pierwszym starcie zaakceptuj warunki; zakładanie konta Docker można \textbf{pominąć} (Skip).
\item Poczekaj, aż ikona wieloryba w zasobniku (przy zegarze) przestanie się animować, a w oknie pojawi się status \textbf{Docker Desktop is running} / ,,Engine running''.
\end{enumerate}

\uwaga{Docker Desktop musi być \textbf{uruchomiony} za każdym razem, gdy ma działać Home Assistant. Warto włączyć autostart: Docker Desktop $\to$ Settings (zębatka) $\to$ General $\to$ \textbf{Start Docker Desktop when you sign in}.}

%=====================================================================
\section{Etap 1 --- Home Assistant w kontenerze}
\label{sec:etap1}

\subsection{Uruchomienie gotowym skryptem}

W folderze projektu \texttt{AMPERE\_POINT} leży skrypt \textbf{\texttt{ha-docker-setup.bat}}. Upewnij się, że Docker Desktop działa, i kliknij skrypt dwukrotnie. Skrypt: sprawdzi Dockera, utworzy folder \texttt{C:\textbackslash HomeAssistant\textbackslash config}, usunie ewentualny stary kontener i uruchomi nowy.

Równoważne polecenie ręczne (Terminal, \emph{bez} administratora), gdyby skryptu zabrakło:

\begin{lstlisting}
mkdir C:\HomeAssistant\config
docker run -d --name homeassistant --restart=unless-stopped -e TZ=Europe/Warsaw -v "C:\HomeAssistant\config:/config" -p 8123:8123 ghcr.io/home-assistant/home-assistant:stable
\end{lstlisting}

Co znaczą kluczowe fragmenty: \texttt{-{}-restart=unless-stopped} --- kontener wstaje sam po restarcie komputera; \texttt{-v "C:/\ldots:/config"} --- konfiguracja HA trafia do zwykłego folderu Windows (przeżywa aktualizacje i przenosiny); \texttt{-p 8123:8123} --- panel HA będzie pod portem 8123.

Pierwszy start pobiera obraz HA (ok.~1,5 GB) --- przez WSL2 zwykle kilka minut. Potem otwórz w przeglądarce:

\begin{center}\Large\texttt{http://localhost:8123}\end{center}

Z innych urządzeń w tej samej sieci: \texttt{http://<IP-komputera>:8123}. Jeżeli strona nie wstaje od razu --- odczekaj 1--2 minuty i odśwież (HA przy pierwszym starcie buduje bazę).

\subsection{Onboarding --- konto i podstawy}

Kreator powitalny HA:
\begin{enumerate}
\item \textbf{Utwórz konto} --- imię, nazwa użytkownika, silne hasło. To konto \textbf{lokalne} (nie wysyłane do żadnej chmury); zapisz dane w menedżerze haseł. Konta i hasła zakłada \textbf{właściciel systemu osobiście} --- nie asystent, nie osoby trzecie.
\item \textbf{Lokalizacja} --- wybierz dom/warsztat (strefa czasowa i wschody/zachody słońca dla automatyzacji), jednostki metryczne.
\item Zgody na statystyki --- wg uznania (można wszystko odznaczyć).
\item Wykryte urządzenia --- na razie \textbf{pomiń} (Zakończ). Dodamy wszystko po kolei świadomie.
\end{enumerate}

\subsection{Jak restartować Home Assistant (w tej instalacji)}
\label{sec:restart}

Wersja kontenerowa HA \textbf{nie ma} przycisku ,,Uruchom ponownie'' w menu System. Restart wykonuje się w \textbf{Docker Desktop}: zakładka \textbf{Containers} $\to$ wiersz \texttt{homeassistant} $\to$ ikona \textbf{Restart} (okrągła strzałka). Alternatywnie w terminalu:

\begin{lstlisting}
docker restart homeassistant
\end{lstlisting}

Po restarcie panel wraca po ok.~30--90 s. Ta czynność będzie potrzebna kilka razy w dalszych etapach.

%=====================================================================
\section{Etap 2 --- HACS (sklep z dodatkami)}
\label{sec:etap2}

\subsection{Instalacja plików HACS}

W folderze projektu jest skrypt \textbf{\texttt{ha-docker-hacs.bat}} --- dwuklik. Skrypt wykonuje wewnątrz kontenera oficjalny instalator HACS i restartuje kontener. Odpowiednik ręczny:

\begin{lstlisting}
docker exec homeassistant bash -c "wget -O - https://get.hacs.xyz | bash -"
docker restart homeassistant
\end{lstlisting}

\subsection{Dodanie integracji HACS i autoryzacja GitHub}

\begin{enumerate}
\item Po restarcie: \textbf{Ustawienia $\to$ Urządzenia i usługi $\to$ Dodaj integrację} $\to$ wyszukaj \textbf{HACS}.
\item Zaakceptuj warunki (checkboxy) $\to$ \textbf{Prześlij}.
\item HA pokaże \textbf{kod urządzenia} i adres \url{https://github.com/login/device}. Otwórz adres, zaloguj się na GitHub, przepisz kod, kliknij \textbf{Authorize}.
\item Wróć do HA $\to$ integracja HACS pojawi się na liście, a w menu bocznym przybędzie pozycja \textbf{HACS}.
\end{enumerate}

\wskazowka{Jeśli po restarcie HACS nie ma na liście integracji --- odśwież stronę przeglądarki twardo (Ctrl+F5). Jeśli nadal brak: sprawdź, czy w \texttt{C:\textbackslash HomeAssistant\textbackslash config\textbackslash custom\_components} istnieje folder \texttt{hacs}, i powtórz restart.}

%=====================================================================
\section{Etap 3 --- oficjalna integracja Tuya}
\label{sec:etap3}

Cel: HA widzi wszystkie ładowarki z konta Smart Life (na razie z jedną encją każda).

\begin{enumerate}
\item \textbf{Ustawienia $\to$ Urządzenia i usługi $\to$ Dodaj integrację} $\to$ \textbf{Tuya}.
\item Kreator poprosi o \textbf{kod użytkownika} (User Code). Znajdziesz go w telefonie: \textbf{Smart Life $\to$ Ja $\to$ zębatka (Ustawienia) $\to$ Konto i bezpieczeństwo $\to$ Kod użytkownika}. Przepisz \textbf{z zachowaniem wielkości liter}.
\item HA wyświetli \textbf{kod QR}. W Smart Life: \textbf{Ja $\to$ ikona skanera} (prawy górny róg) $\to$ zeskanuj ekran $\to$ \textbf{zatwierdź logowanie} na telefonie.
\item Po chwili integracja pokaże urządzenia konta (u nas 10: Q11/Q22/Q37/GPEV100 itd.).
\end{enumerate}

\uwaga{Każde urządzenie ma teraz tylko \textbf{jedną encję} (przełącznik). \textbf{To jest oczekiwane.} Brakujące czujniki doda xtend\_tuya w etapach 4--7. Nie ma sensu szukać błędu ani restartować --- tak działa oficjalna integracja z tymi ładowarkami.}

%=====================================================================
\section{Etap 4 --- xtend\_tuya, część A (bez chmury deweloperskiej)}
\label{sec:etap4}

Cel: encje rosną z 1 do kilkunastu--dwudziestu kilku na urządzenie (u nas łącznie z 10 do ok.~154).

\subsection{Pobranie w HACS}

\begin{enumerate}
\item Menu boczne $\to$ \textbf{HACS} $\to$ wyszukaj \textbf{Xtend Tuya}.
\item Otwórz pozycję $\to$ przycisk \textbf{Pobierz / Download} (prawy dolny róg) $\to$ potwierdź.
\item \textbf{Zrestartuj kontener} (rozdz.~\ref{sec:restart}).
\end{enumerate}

\uwaga{Ekran pozycji w HACS pokazuje README z przyciskiem ,,OPEN HACS REPOSITORY'', który \textbf{niczego nie instaluje} --- to tylko odnośnik. Po pobraniu przycisk ,,Download'' \textbf{znika} --- to nie awaria, tak wygląda stan ,,pobrane''. Pobranie w HACS kopiuje wyłącznie pliki na dysk; encje pojawią się dopiero po \textbf{dodaniu integracji} (poniżej). To rozróżnienie było źródłem wielotygodniowej frustracji --- pliki leżały na dysku, a integracji nikt nie ,,dokończył''.}

\subsection{Dodanie integracji --- z kluczowym oknem, którego nie wolno zamknąć}
\label{sec:krytyczne-okno}

\begin{enumerate}
\item \textbf{Ustawienia $\to$ Urządzenia i usługi $\to$ Dodaj integrację} $\to$ wyszukaj ,,xtend'' $\to$ wybierz \textbf{Xtend Tuya} (nie mylić z ,,TuyaExtend AmperePoint (helper)'' --- ten będzie w etapie 8).
\item Podaj \textbf{kod użytkownika} (ten sam co w etapie 3) $\to$ zeskanuj \textbf{kod QR} aplikacją Smart Life $\to$ zatwierdź na telefonie.
\item \textbf{MOMENT KRYTYCZNY.} Pojawi się formularz \textbf{,,Add Tuya OpenAPI credentials''}. Wygląda jak opcjonalny i ma krzyżyk --- ale \textbf{zamknięcie go krzyżykiem po cichu anuluje cały kreator}: wpis integracji nie powstaje, HA nie pokazuje żadnego błędu, a użytkownik jest przekonany, że ,,wszystko zrobił''. Dokładnie to zjawisko zablokowało wcześniejsze próby (odtworzone i potwierdzone w sesji 2026-07-05).
\item W formularzu (skoro kluczy deweloperskich jeszcze nie mamy):
  \begin{itemize}
  \item zaznacz \textbf{,,Don't ask for OpenAPI credentials''},
  \item \textbf{mimo zaznaczenia} wybierz \textbf{Country = Poland} oraz \textbf{endpoint = Europe} (\texttt{https://openapi.tuyaeu.com}) --- formularz ma błąd walidacji i bez tych pól zgłosi niezrozumiały komunikat ,,value must be one of [\ldots]'',
  \item zatwierdź.
  \end{itemize}
\item Okno \textbf{,,Name and assign''} (nazwy/obszary) --- zostaw domyślne $\to$ \textbf{Zakończ}.
\end{enumerate}

\subsection{Weryfikacja etapu A}

\textbf{Ustawienia $\to$ Urządzenia i usługi}: na liście musi być wpis \textbf{Xtend Tuya --- 10 urządzeń}. Liczby encji po etapie A z wdrożenia wzorcowego:

\begin{center}
\begin{tabular}{lcc}
\toprule
\textbf{Urządzenie} & \textbf{encje przed} & \textbf{encje po etapie A}\\
\midrule
Q22 OTA & 1 & 24\\
Q22 OTA 2 & 1 & 24\\
urzadzenie\_zepsute (Q22 OTA) & 1 & 24\\
GPEV100 & 1 & 16\\
Q11 OTA & 1 & 15\\
ladowarka\_q11 (Q11 OTA) & 1 & 15\\
Q37 OTA & 1 & 15\\
EV Charger VE & 1 & 15\\
EV Charging Station & 1 & 5\\
AmperePoint EV Charger & 1 & 1 (patrz rozdz.~\ref{sec:lightsource})\\
\bottomrule
\end{tabular}
\end{center}

Jeśli wpisu Xtend Tuya \textbf{nie ma} --- kreator został przerwany (najpewniej na formularzu OpenAPI); powtórz podrozdział~\ref{sec:krytyczne-okno} od początku.

%=====================================================================
\section{Etap 5 --- projekt na platformie deweloperskiej Tuya}
\label{sec:etap5}

Cel: drugi, pełniejszy kanał dostępu do chmury (OpenAPI) --- klucze \textbf{Client ID + Secret}, dzięki którym xtend\_tuya zobaczy ,,prawie wszystkie'' punkty danych.

\begin{enumerate}
\item Wejdź na \url{https://iot.tuya.com} i \textbf{załóż darmowe konto} (może być ten sam e-mail co w Smart Life; po zalogowaniu serwis przekierowuje na \texttt{platform.tuya.com}).
\item \textbf{Cloud $\to$ Development $\to$ Create Cloud Project}. Wypełnij:
  \begin{itemize}
  \item Project Name: np. \textbf{AmperePoint HA};
  \item Industry: \textbf{Smart Home}; Development Method: \textbf{Smart Home};
  \item Data Center: \textbf{Central Europe Data Center} --- wybór \textbf{krytyczny}: musi odpowiadać regionowi konta Smart Life (dla Polski --- Europa Środkowa). Darmowy plan pozwala na jedno centrum danych.
  \end{itemize}
\item Kreator \textbf{Authorize API Services} zaproponuje usługi --- zatwierdź proponowane bez zmian (IoT Core, Authorization Token Management, Smart Home Basic Service, Data Dashboard Service, ewentualnie Smart Home Scene Linkage oznaczona jako przestarzała --- nieszkodliwa).
\item Po utworzeniu projektu, karta \textbf{Overview}, sekcja \emph{Authorization Key}: \textbf{Access ID/Client ID} oraz \textbf{Access Secret/Client Secret} (odsłania ikona oka). \uwaga{Secret pełni rolę hasła do API: nie wysyłać nikomu, nie wklejać do dokumentów współdzielonych, trzymać w menedżerze haseł. W razie wycieku: Overview $\to$ ikona przy kluczu $\to$ reset sekretu, potem aktualizacja w HA.}
\item Powiąż konto aplikacji: zakładka \textbf{Devices $\to$ Link App Account $\to$ Add App Account $\to$ Tuya App Account Authorization} $\to$ zeskanuj kod QR aplikacją Smart Life $\to$ zatwierdź. Wybierz tryb \textbf{Automatic Link} (dołącza też urządzenia dodane w przyszłości). Platforma potwierdzi listę urządzeń konta.
\end{enumerate}

%=====================================================================
\section{Etap 6 --- tryb DP Instruction (pełne punkty danych)}
\label{sec:etap6}

Chmura Tuya wystawia urządzenia przez API w trybie \textbf{Standard Instruction} (ujednolicone, ale okrojone) albo \textbf{DP Instruction} (surowe punkty danych --- pełna lista zmiennych urządzenia). Dokumentacja xtend\_tuya zaleca DP. Procedura --- \textbf{dla każdej kategorii produktu osobno}:

\begin{enumerate}
\item Na platformie: projekt $\to$ \textbf{Devices $\to$ All Devices} $\to$ przycisk \textbf{,,\ldots''} przy pasku \emph{View Devices by Product}.
\item Najedź kursorem na kafelek produktu $\to$ ikona \textbf{ołówka}.
\item Zaznacz \textbf{DP Instruction} $\to$ \textbf{Save Configuration} $\to$ \textbf{OK}.
\item Powtórz dla wszystkich produktów z listy (u nas: EV Charger\_VE, AmperePoint EV Charger, Q20\_3.5KW\_with\_OTA, EV Charging Station, Q11\_11KW\_with\_OTA, Q22 OTA, GPEV100).
\end{enumerate}

Ostrzeżenie platformy (,,devices might not be able to be controlled by the original standard instructions'') dotyczy \textbf{wyłącznie tego projektu API} --- aplikacja Smart Life i same ładowarki działają bez zmian.

\subsection{Znany przypadek: produkt zarejestrowany jako źródło światła}
\label{sec:lightsource}

Produkt \textbf{AmperePoint EV Charger} (PID \texttt{gbmxngploofmhbjc}) jest w chmurze Tuya zarejestrowany przez producenta w kategorii \textbf{,,Light Source''} z jednym jedynym punktem danych \texttt{switch\_led}. Dlatego to urządzenie ma i będzie miało \textbf{tylko 1 encję}, niezależnie od wszystkich kroków po stronie HA. Rozwiązanie leży wyłącznie po stronie producenta: zgłosić prośbę o przeniesienie produktu do kategorii ładowarek EV z pełnym zestawem DP (podając PID).

%=====================================================================
\section{Etap 7 --- klucze OpenAPI w HA i weryfikacja końcowa}
\label{sec:etap7}

\begin{enumerate}
\item \textbf{Ustawienia $\to$ Urządzenia i usługi $\to$ Xtend Tuya $\to$ Konfiguruj}. Otworzy się znany formularz ,,Add Tuya OpenAPI credentials''.
\item Tym razem: \textbf{odznacz} ,,Don't ask for OpenAPI credentials''; Country = \textbf{Poland}; endpoint = \textbf{Europe}; \textbf{Tuya IoT Access ID} i \textbf{Access Secret} wklej z karty Overview projektu (etap 5). Zatwierdź.
\item \textbf{Zrestartuj kontener} (rozdz.~\ref{sec:restart}) i odczekaj 1--2 minuty.
\item Weryfikacja: \textbf{Ustawienia $\to$ Urządzenia i usługi $\to$ Xtend Tuya} --- liczby encji przy urządzeniach powinny wzrosnąć (Q22 docelowo 30+). W razie braków: odśwież stronę; potem menu \textbf{,,\ldots''} przy wpisie Xtend Tuya $\to$ \textbf{Załaduj ponownie}; na końcu zajrzyj do logów (\textbf{Ustawienia $\to$ System $\to$ Logi}, filtr ,,xtend'') --- typowe błędy w tabeli rozdz.~\ref{sec:problemy}.
\end{enumerate}

%=====================================================================
\section{Etap 8 --- dodatek firmowy TuyaExtend AmperePoint}
\label{sec:etap8}

Dodatek firmowy \textbf{nie pobiera} danych z Tuya --- mapuje istniejące encje i dokłada warstwę prezentacji: koszty sesji wg taryfy, statusy po polsku, jedno ,,ładne'' urządzenie ładowarki. Konfigurować \textbf{dopiero po} etapie 7, gdy encje źródłowe istnieją.

\subsection{Instalacja plików (jeśli jeszcze nie ma)}

Repozytorium firmowe jest \textbf{prywatne}, więc HACS go nie zobaczy (HACS czyta tylko repozytoria publiczne --- stąd wcześniejsze błędy 404 przy próbach dodania adresów \texttt{amperepoint/tuya-ap} i podobnych). Instalacja ręczna: skopiuj folder \texttt{tuyaextend\_amperepoint} (z folderu projektu AMPERE\_POINT) do \texttt{C:\textbackslash HomeAssistant\textbackslash config\textbackslash custom\_components\textbackslash} tak, by powstała ścieżka \texttt{custom\_components\textbackslash tuyaextend\_amperepoint\textbackslash manifest.json}. Zrestartuj kontener.

\subsection{Kreator konfiguracji}

\textbf{Dodaj integrację $\to$ TuyaExtend AmperePoint}, następnie:

\begin{enumerate}
\item \textbf{Model}: np. AmperePoint Q22 OTA; \textbf{nazwa} urządzenia (np. ,,Ładowarka Q22 --- stanowisko 1'').
\item \textbf{Taryfa}: cena prądu w \textbf{PLN/kWh}; waluta \textbf{PLN}.
\item Parametry pozostawić domyślne: próg mocy ładowania \textbf{0{,}25 kW}, minuty bezczynności \textbf{3}, tryb energii sesji \textbf{Auto}.
\item \textbf{Mapowanie encji źródłowych}: przełącznik ładowania, moc, energia sesji, energia całkowita, limit prądu, status, błąd, napięcia i prądy L1--L3 --- wybieraj encje utworzone przez xtend\_tuya dla danej ładowarki.
\end{enumerate}

Pomyłkę w mapowaniu naprawia się przyciskiem \textbf{Konfiguruj} przy wpisie integracji --- bez usuwania czegokolwiek.

%=====================================================================
\section{Etap 9 --- porządki, panel Energia, pulpit}
\label{sec:etap9}

\subsection{Rejestr encji --- nazwy, duplikaty, obszary}

\textbf{Ustawienia $\to$ Urządzenia i usługi $\to$ zakładka Encje}:
\begin{itemize}
\item nadaj kluczowym encjom czytelne nazwy i identyfikatory (np. \texttt{sensor.q22\_moc});
\item \textbf{duplikaty}: oficjalna Tuya i xtend\_tuya potrafią wystawić ten sam przełącznik dwa razy --- to normalne; zbędną encję \textbf{wyłącz} (nie usuwaj integracji);
\item wyłącz encje niepotrzebne (mniej szumu w bazie danych);
\item przypisz urządzenia do \textbf{obszarów} (np. ,,Pole'', ,,Serwis'').
\end{itemize}

\subsection{Panel Energia}

\textbf{Ustawienia $\to$ Pulpity $\to$ Energia} $\to$ sekcja ,,Poszczególne urządzenia'' $\to$ dodaj encje energii ładowarek (sumujące kWh --- zwykle ,,energia całkowita'' / \texttt{forward\_energy\_total}). HA zacznie budować dzienne i miesięczne wykresy --- podstawa analiz i rozliczeń. Wykresy pojawiają się po 1--2 godzinach zbierania danych.

\subsection{Pulpit ładowarek}

Na pulpicie (przycisk \textbf{Edytuj} w prawym górnym rogu) warto zestawić: kartę \textbf{wskaźnika} (gauge) z mocą chwilową, kartę \textbf{historii} z energią sesji, kartę \textbf{encji} ze statusem, limitem prądu i przełącznikiem ładowania.

%=====================================================================
\section{Etap 10 --- automatyzacje}
\label{sec:etap10}

Automatyzacja składa się z \textbf{wyzwalacza} (,,kiedy?''), opcjonalnych \textbf{warunków} (,,czy na pewno?'') i \textbf{akcji} (,,co zrobić?''). Buduje się je w kreatorze: \textbf{Ustawienia $\to$ Automatyzacje i sceny $\to$ Utwórz automatyzację}. Poniżej trzy sprawdzone przykłady w zapisie YAML (kreator ma tryb ,,Edytuj w YAML''). \textbf{Identyfikatory encji podmień na własne} z Rejestru encji. Powiadomienia wymagają aplikacji \textbf{Home Assistant Companion} na telefonie.

\textbf{1) Powiadomienie o zakończeniu ładowania:}

\begin{lstlisting}
alias: Powiadom o koncu ladowania Q22
trigger:
  - platform: state
    entity_id: sensor.q22_status
    from: "charging"
    to: "charger_free"
action:
  - service: notify.mobile_app_telefon_dawida
    data:
      title: "Ladowarka Q22"
      message: "Ladowanie zakonczone."
\end{lstlisting}

\textbf{2) Ładowanie tylko w taniej taryfie (22:00--6:00):}

\begin{lstlisting}
alias: Q22 tania taryfa
trigger:
  - platform: time
    at: "22:00:00"
    id: start
  - platform: time
    at: "06:00:00"
    id: stop
action:
  - choose:
      - conditions: [{ condition: trigger, id: start }]
        sequence:
          - service: switch.turn_on
            target: { entity_id: switch.q22_przelacznik }
      - conditions: [{ condition: trigger, id: stop }]
        sequence:
          - service: switch.turn_off
            target: { entity_id: switch.q22_przelacznik }
\end{lstlisting}

\textbf{3) Alarm przy błędzie ładowarki:}

\begin{lstlisting}
alias: Alarm bledu ladowarki
trigger:
  - platform: state
    entity_id: sensor.q22_blad
condition:
  - condition: template
    value_template: "{{ trigger.to_state.state not in ['0','ok','unknown','unavailable'] }}"
action:
  - service: notify.mobile_app_telefon_dawida
    data:
      title: "UWAGA: blad ladowarki"
      message: "Kod bledu: {{ trigger.to_state.state }}"
\end{lstlisting}

Kierunki rozwoju: obniżanie limitu prądu w szczycie (encje typu \emph{number}, usługa \texttt{number.set\_value}); dzienny raport energii; \textbf{ładowanie nadwyżką z fotowoltaiki} --- wyzwalaczem jest moc eksportu do sieci (patrz symulator, rozdz.~\ref{sec:pvsim}).

%=====================================================================
\section{Rozszerzenie opcjonalne: kamera Imou IPC-T42EP}
\label{sec:kamera}

Zidentyfikowany sprzęt: \textbf{Imou Turret SE 4MP, model IPC-T42EP}, obiektyw 2,8 mm, zasilacz 12 V/0,5 A. Kamera ma \textbf{port Ethernet} i Wi-Fi 2,4 GHz, \textbf{oficjalne wsparcie ONVIF}, slot microSD (do 256 GB) i mikrofon --- to jeden z łatwiejszych przypadków integracyjnych.

\uwaga{Stan na 2026-07-05: egzemplarz jest \textbf{powiązany z cudzym kontem Imou} i nie da się go dodać do własnego. Zanim wykonasz kroki poniżej: poprzedni właściciel musi usunąć kamerę ze swojego konta w aplikacji Imou Life (Ustawienia urządzenia $\to$ Usuń), \emph{albo} wykonaj reset fabryczny przyciskiem reset przy porcie kabla ($\sim$10 s do sygnału) i sparuj ponownie; jeśli blokada konta pozostanie --- konieczny kontakt ze wsparciem Imou z dowodem zakupu.}

Ścieżka zalecana (lokalna, podgląd bez chmury):
\begin{enumerate}
\item Sparuj kamerę w aplikacji \textbf{Imou Life} (najlepiej po \textbf{kablu LAN}; przy Wi-Fi pamiętaj o 2,4 GHz).
\item W routerze nadaj kamerze \textbf{stałe IP}.
\item W Imou Life sprawdź ustawienia kamery i \textbf{wyłącz szyfrowanie/TLS}, jeśli występuje (nowszy firmware blokuje zwykły RTSP przy włączonym szyfrowaniu).
\item W HA: \textbf{Dodaj integrację $\to$ ONVIF} $\to$ IP kamery, port 80, login \texttt{admin}, hasło = \textbf{safety code z naklejki} na spodzie kamery lub hasło ustawione w aplikacji.
\item Jeśli ONVIF zawiedzie --- \textbf{Generic Camera} ze strumieniem RTSP:
\begin{lstlisting}
rtsp://admin:HASLO@IP_KAMERY:554/cam/realmonitor?channel=1&subtype=0
\end{lstlisting}
(\texttt{subtype=1} = lżejszy strumień pomocniczy).
\end{enumerate}

Ścieżka dodatkowa (chmurowa --- przełączniki, detekcja ruchu jako encje): darmowe konto na \url{https://open.imoulife.com} $\to$ \textbf{My App} $\to$ \textbf{AppId + AppSecret} (analogia 1:1 do kluczy Tuya) $\to$ w HACS integracja \textbf{Imou Life}.

%=====================================================================
\section{Rozszerzenie opcjonalne: symulator falownika SUN2000}
\label{sec:pvsim}

Do ćwiczeń z fotowoltaiką \textbf{przed} zakupem falownika powstał symulator \textbf{SUN2000-10KTL-M1}: serwer Modbus TCP w czystym Pythonie (zero bibliotek do instalowania), pliki w folderze projektu \texttt{pv\_sim}. HA łączy się z nim \textbf{prawdziwą} integracją Huawei Solar --- dokładnie tak, jak z fizycznym falownikiem.

\begin{enumerate}
\item Uruchomienie (terminal):
\begin{lstlisting}
cd C:\Users\Lenovo\Desktop\AMPERE_POINT\pv_sim
python sun2000_sim.py
\end{lstlisting}
Gdy zapora Windows zapyta --- \textbf{Zezwól} (sieci prywatne). Przydatne parametry: \texttt{-{}-kwp 10} (moc instalacji), \texttt{-{}-time-scale 60} (doba mija w 24 min), \texttt{-{}-port 5020} (gdy 502 zajęty), \texttt{-{}-no-meter} (bez licznika).
\item W HACS pobierz \textbf{Huawei Solar} (autor wlcrs) $\to$ restart kontenera.
\item \textbf{Dodaj integrację $\to$ Huawei Solar} $\to$ typ \textbf{Modbus TCP}: host \texttt{host.docker.internal} (tak kontener HA ,,widzi'' Windowsa), port \texttt{502}, slave/device ID \texttt{0} lub \texttt{1}, tryb bez podwyższonych uprawnień.
\item Po chwili pojawi się urządzenie SUN2000 z encjami: moc czynna, moc DC, energie, fazy L1--L3, stringi PV, temperatura, status oraz licznik (eksport/import). Symulator generuje realistyczną krzywą dobową z ,,chmurami''; energie liczników przeżywają restart (plik \texttt{sun2000\_sim\_state.json}).
\item Test bez HA: \texttt{python test\_client.py 127.0.0.1 502}. Szczegóły i szkic automatyzacji ,,ładowanie nadwyżką'': \texttt{pv\_sim\textbackslash README\_URUCHOMIENIE.md}.
\end{enumerate}

Po zakupie prawdziwego falownika wystarczy w integracji podmienić host na IP falownika/SDongle --- reszta konfiguracji (encje, panel Energia, automatyzacje) zostaje.

%=====================================================================
\section{Kopie zapasowe i aktualizacje}
\label{sec:backup}

\begin{description}[style=nextline]
\item[Kopia zapasowa] najpewniejsza metoda w instalacji kontenerowej: zatrzymać kontener (Docker Desktop $\to$ Stop) i \textbf{skopiować cały folder} \texttt{C:\textbackslash HomeAssistant\textbackslash config} (np. na dysk zewnętrzny), potem Start. Kopię robić przed każdą większą zmianą (aktualizacja HA, nowe integracje).
\item[Aktualizacja Home Assistant] w terminalu:
\begin{lstlisting}
docker pull ghcr.io/home-assistant/home-assistant:stable
docker rm -f homeassistant
\end{lstlisting}
\ldots{}a następnie ponownie polecenie \texttt{docker run} z rozdz.~\ref{sec:etap1} (lub \texttt{ha-docker-setup.bat} --- UWAGA: skrypt usuwa stary kontener, ale konfiguracja w \texttt{C:\textbackslash HomeAssistant\textbackslash config} zostaje nietknięta).
\item[Aktualizacje xtend\_tuya] przez HACS, ale \textbf{rozważnie}: dodatek korzysta z nieoficjalnych mechanizmów i bywa czuły na zmiany HA. Po każdej większej aktualizacji HA sprawdzić w HACS zgodność wersji; w razie awarii --- tymczasowo wrócić do poprzedniego obrazu HA.
\item[Przypomnienie okresowe] darmowy plan \textbf{IoT Core Trial Edition} na platformie Tuya wygasa (typowo co pół roku). Wpisz do kalendarza: \textbf{Cloud $\to$ projekt $\to$ IoT Core $\to$ Extend Trial} (bezpłatne).
\end{description}

%=====================================================================
\section{Rozwiązywanie problemów}
\label{sec:problemy}

\begin{longtable}{>{\raggedright\arraybackslash}p{4.2cm} >{\raggedright\arraybackslash}p{4.6cm} >{\raggedright\arraybackslash}p{5.6cm}}
\toprule
\textbf{Objaw} & \textbf{Prawdopodobna przyczyna} & \textbf{Rozwiązanie}\\
\midrule
\endhead
\multicolumn{3}{l}{\textit{--- Instalacja podstawy (etapy 0--2) ---}}\\
\addlinespace
\texttt{docker} ,,nie jest rozpoznawalny'' / skrypt zgłasza błąd & Docker Desktop nie działa albo silnik jeszcze startuje & Uruchomić Docker Desktop, poczekać na ,,running'', powtórzyć. \\
\addlinespace
\texttt{localhost:8123} nie odpowiada & Kontener startuje / pierwszy start buduje bazę & Odczekać 1--2 min; Docker Desktop $\to$ Containers $\to$ status \texttt{homeassistant} ma być ,,Running''; logi kontenera przyciskiem ,,View details''.\\
\addlinespace
Port 8123 zajęty (błąd przy \texttt{docker run}) & Inny program używa portu & W poleceniu zmienić mapowanie na \texttt{-p 8124:8123} i wchodzić na \texttt{localhost:8124}.\\
\addlinespace
HACS nie pojawia się po instalacji & Restart wykonany przed końcem skryptu; twardy cache przeglądarki & Sprawdzić folder \texttt{custom\_components\textbackslash hacs}; \texttt{docker restart homeassistant}; Ctrl+F5.\\
\addlinespace
\multicolumn{3}{l}{\textit{--- Tuya / xtend\_tuya (etapy 3--7) ---}}\\
\addlinespace
Po ,,pobraniu'' xtend\_tuya encji nie przybywa & Pobranie $\ne$ dodanie: brak wpisu integracji (kreator przerwany, zwykle krzyżykiem na formularzu OpenAPI) & Przejść kreator do końca wg rozdz.~\ref{sec:krytyczne-okno}; restart nic nie da, dopóki wpis nie istnieje.\\
\addlinespace
Formularz OpenAPI z błędem ,,value must be one of [\ldots]'' & Błąd walidacji: kraj/endpoint wymagane mimo checkboxa & Wybrać Country=Poland i endpoint=Europe, dopiero wtedy zatwierdzić.\\
\addlinespace
Błąd ,,sign invalid'' w logach & Zły Client ID/Secret (literówka, spacja) albo zły endpoint & Skopiować klucze ponownie z Overview; endpoint = Europe (\texttt{openapi.tuyaeu.com}).\\
\addlinespace
Błąd ,,permission deny'' / kod 1106 & Nieautoryzowana usługa API albo wygasły trial & Platforma $\to$ Service API: IoT Core i Authorization aktywne; w razie potrzeby Authorize/Renew.\\
\addlinespace
Encje przestają się aktualizować po miesiącach & Wygasł IoT Core Trial Edition & Cloud $\to$ projekt $\to$ IoT Core $\to$ Extend Trial; wpisać przypomnienie do kalendarza.\\
\addlinespace
Na platformie ,,No data found'' / puste listy & Przeglądarka zgubiła region (przekierowania platform $\leftrightarrow$ eu.platform) & Odświeżyć; sprawdzić wybór Central Europe DC w prawym górnym rogu.\\
\addlinespace
QR z platformy nie działa / błąd autoryzacji & Konto aplikacji w innym centrum danych niż projekt & Projekt musi mieć ten sam DC co konto Smart Life (dla PL: Central Europe).\\
\addlinespace
Wszystkie encje ,,niedostępne'' & Brak internetu / awaria chmury / urządzenie offline & Sprawdzić internet i stan w Smart Life; integracja jest chmurowa.\\
\addlinespace
Po aktualizacji HA Xtend Tuya nie startuje & Niezgodność wersji dodatku z HA & HACS $\to$ aktualizacja xtend\_tuya; brak poprawki $\to$ powrót do poprzedniego obrazu HA.\\
\addlinespace
Zdublowane encje (dwa przełączniki) & Tuya i xtend\_tuya wystawiają ten sam element & Normalne; zbędną encję wyłączyć w Rejestrze encji.\\
\addlinespace
,,AmperePoint EV Charger'' ma wciąż 1 encję & Produkt zarejestrowany jako Light Source (1 DP) & Zgłosić producentowi (PID \texttt{gbmxngploofmhbjc}); po stronie HA nie ma co naprawiać.\\
\addlinespace
Ładowarka zniknęła po zmianie routera & Moduły Tuya wymagają 2,4 GHz i ponownego parowania & Włączyć 2,4 GHz, sparować w Smart Life; wróci do HA sama (Automatic Link).\\
\addlinespace
\multicolumn{3}{l}{\textit{--- Rozszerzenia ---}}\\
\addlinespace
Kamera Imou: ,,urządzenie już dodane'' przy parowaniu & Egzemplarz powiązany z cudzym kontem & Usunięcie z poprzedniego konta albo reset + parowanie; ostatecznie wsparcie Imou (dowód zakupu).\\
\addlinespace
Kamera: RTSP nie łączy & Włączone szyfrowanie/TLS w Imou Life; złe hasło & Wyłączyć szyfrowanie; hasło = safety code z naklejki lub z aplikacji.\\
\addlinespace
Symulator: HA nie łączy się & Zapora Windows; zły host; port zajęty & Zezwolić Pythonowi (sieci prywatne); host = \texttt{host.docker.internal}; ew. \texttt{-{}-port 5020} i ten sam port w HA.\\
\bottomrule
\end{longtable}

%=====================================================================
\section{Checklista końcowa --- ,,czy wszystko działa''}

\begin{center}
\begin{tabular}{>{\raggedright\arraybackslash}p{11.2cm} c}
\toprule
\textbf{Punkt kontrolny} & \textbf{OK?}\\
\midrule
Docker Desktop startuje z systemem; kontener \texttt{homeassistant} = Running & $\square$\\
\texttt{http://localhost:8123} loguje do panelu HA & $\square$\\
HACS widoczny w menu bocznym & $\square$\\
Integracja Tuya: wszystkie ładowarki na liście urządzeń & $\square$\\
Wpis \textbf{Xtend Tuya} istnieje na liście integracji (nie tylko ,,pobrane'' w HACS) & $\square$\\
Q22: 30+ encji; pozostałe ładowarki: kilkanaście+ & $\square$\\
Klucze OpenAPI wpisane; w logach brak ,,sign invalid'' / ,,permission deny'' & $\square$\\
Wszystkie produkty na platformie w trybie \textbf{DP Instruction} & $\square$\\
TuyaExtend AmperePoint skonfigurowany; koszt sesji liczy się w PLN & $\square$\\
Panel Energia pokazuje energię ładowarek (po 1--2 h) & $\square$\\
Testowa automatyzacja (np. powiadomienie) zadziałała & $\square$\\
Kopia zapasowa \texttt{C:\textbackslash HomeAssistant\textbackslash config} wykonana & $\square$\\
Przypomnienie ,,Extend Trial'' (IoT Core) w kalendarzu & $\square$\\
\bottomrule
\end{tabular}
\end{center}

%=====================================================================
\section{Słowniczek pojęć i skrótów}
\label{sec:slowniczek}

\begin{description}[style=nextline,itemsep=1pt]
\item[API] interfejs programistyczny --- ,,okienko podawcze'', przez które programy proszą usługę o dane lub działania.
\item[Centrum danych (DC)] serwerownia regionu przypisana do konta Tuya/Imou; u nas Central Europe.
\item[Client ID / Client Secret] para ,,login + hasło'' projektu deweloperskiego do podpisywania żądań API (Tuya); odpowiednik w Imou: AppId/AppSecret.
\item[Config entry (wpis konfiguracyjny)] zapisany wynik kreatora integracji w HA; \textbf{bez niego integracja nie działa}, choćby pliki leżały na dysku.
\item[Docker / kontener / obraz] kontener = odizolowane ,,pudełko'' z programem; obraz = wzorzec, z którego kontener powstaje; Docker Desktop = program zarządzający tym na Windows.
\item[DP (data point)] pojedyncza zmienna urządzenia Tuya (moc, status\ldots); tryb DP Instruction wystawia pełną listę DP przez API.
\item[Encja (entity)] pojedynczy czujnik/przełącznik/wartość w HA.
\item[Endpoint] adres serwera API (u nas \texttt{https://openapi.tuyaeu.com}).
\item[HA / HAOS] Home Assistant; HAOS = jego samodzielny system operacyjny (używany w wariancie VM, u nas zastąpiony kontenerem).
\item[HACS] społecznościowy instalator dodatków do HA.
\item[Integracja] moduł HA łączący go z usługą lub urządzeniami.
\item[Modbus TCP] prosty protokół odczytu/zapisu rejestrów po sieci; tak rozmawia się z falownikami (i naszym symulatorem).
\item[ONVIF] standard komunikacji kamer IP; pozwala dodać kamerę do HA bez aplikacji producenta.
\item[OpenAPI (Tuya)] oficjalne API deweloperskie chmury Tuya.
\item[PID / Device ID] identyfikator modelu produktu / konkretnego egzemplarza w Tuya.
\item[QR] kod skanowany telefonem; tu: autoryzacja logowań Tuya/Imou.
\item[RTSP] protokół strumienia wideo z kamer; adres zaczyna się od \texttt{rtsp://}.
\item[Safety code] fabryczne hasło urządzenia Imou, wydrukowane na naklejce.
\item[UAC] okno Windows ,,Czy chcesz zezwolić tej aplikacji\ldots'' --- prośba o prawa administratora.
\item[User Code (kod użytkownika)] identyfikator konta Smart Life używany przy logowaniu integracji przez QR.
\item[WSL2] lekka maszyna wirtualna z Linuksem, w której Docker Desktop uruchamia kontenery.
\item[YAML] tekstowy format konfiguracji automatyzacji HA (wcięcia mają znaczenie).
\end{description}

\vspace{8pt}
\begin{center}
\textcolor{navy}{\rule{0.6\textwidth}{0.4pt}}\\[4pt]
{\small Koniec dokumentu --- wersja 1, 2026-07-06. Kolejne wersje oznaczać v2, v3\ldots{} bez usuwania poprzednich.\\
Dokumenty powiązane: \texttt{AMPERE\_POINT\_home\_assistant\_xtend\_tuya\_raport\_v1} (wyjaśnienie działania),\\ \texttt{INSTRUKCJA\_HomeAssistant\_Docker.md}, \texttt{pv\_sim\textbackslash README\_URUCHOMIENIE.md}, notatki sesji w \texttt{home\_assistant}.}
\end{center}

\end{document}
